% ---------- פרק 3: סקירת הספרות ומאמר הדגל ----------
% 22.11.2026: גוף שלישי וקיצור. מבנה תמטי וביקורתי, לפי פרק 5 במדריך.
% כל מקור נבדק מול DOI/arXiv או טקסט מלא. סטטוס: "מקורות - רשימה מאומתת.md".
% מאמר הדגל (Sadeghi 2026) והמאמר הנלווה (Zhao 2025) נקראו בטקסט מלא; כל המספרים בסעיפים שלהם לקוחים מהם.
\chapter{סקירת הספרות ומאמר הדגל}\label{ch:literature}

הסקירה מחולקת לפי תחומי ידע, וכל תחום מתאים לחלק אחר בבעיה (פרק~\ref{ch:problem}): הערכה רב-קריטריונית וקביעת משקלות, איכות ספקים, סיכון גיאופוליטי, כשל פיננסי, זיהוי אירועים שליליים, אבטחת סייבר ומודלי שפה בשרשרת אספקה. בסוף הפרק מוצגות מטריצת הספרות, מאמר הדגל, המאמר הנלווה ופער הידע.

% =====================================================================
\section{הערכת ספקים וקבלת החלטות רב-קריטריונית}\label{sec:lit-mcdm}

\en{Ho}, \en{Xu} ו-\en{Dey} סקרו את שיטות קבלת ההחלטות הרב-קריטריוניות (\en{MCDM}) לבחירת ספקים, והסיקו שהן עדיפות על בחירה לפי עלות בלבד \cite{ho2010}. לכן ההמלצה המוצעת נשענת על כמה ממדים בעלי משקלות. השיטה הנפוצה לקביעת משקלות היא תהליך ההיררכיה האנליטית (\en{AHP}) של \en{Saaty} \cite{saaty1990}: ההחלטה מפורקת להיררכיה של מטרה, קריטריונים ותתי-קריטריונים, כל זוג קריטריונים מושווה, והמשקלות מחושבים מההשוואות, תוך בדיקת \textbf{עקביות} השיפוטים. המבנה ההיררכי מתאים למודל המשקלות המוצע (חלק, ממד וקטגוריות, פרק~\ref{ch:design}). \en{Boukrouh} ועמיתיו השוו שלוש דרכים לקביעת משקלות לבחירת ספקים: שיפוט (\en{AHP}), חישוב מהנתונים (אנטרופיה) ושילוב של השתיים \cite{boukrouh2024}. בניתוח רגישות על 21 תרחישים, השילוב נשאר יציב ב-19, החישוב ב-3, והשיפוט בלבד באף אחד.

\textbf{ביקורת.} לשיטות הקלאסיות שתי הנחות שלא מתקיימות כאן. הראשונה היא \textbf{מידע מלא}: השיטות מניחות ציון לכל ספק בכל קריטריון, ואינן אומרות מה לעשות כשלחברה פרטית אין דוחות כספיים או כיסוי בחדשות; בזה עוסקת שאלת המשנה השלישית. השנייה נוגעת ל\textbf{קביעת המשקלות}: בפרויקט זה את ההשוואות בזוגות ממלאות מגישות הדוח, כל השוואה מנומקת במקור מהספרות, והמשקלות מוצגים למנחה התעשייתי להערות לפני שהם נקבעים (\en{T1.3}). כאן נדרשת הכרעה מול ממצא \en{Boukrouh} ועמיתיו: הם ממליצים לשלב שיפוט עם חישוב מהנתונים, אבל שיטת האנטרופיה מחשבת משקלות מפיזור הציונים של \textbf{כמה ספקים} באותו קריטריון, ובשלב טרום ההתקשרות מוערך ספק אחד שאין לו נתוני ביצוע. לכן המשקלות נקבעים בשיפוט, ובגלל הרגישות שהם מצאו בשיפוט בלבד, כל משקל נבדק בניתוח רגישות. אם בהמשך יצטבר מאגר הערכות של ספקים רבים, אפשר יהיה לעבור לשיטה המשולבת.

% =====================================================================
\section{הערכת איכות וביצועים של ספקים}\label{sec:lit-quality}

\en{Dey}, \en{Bhattacharya} ו-\en{Ho} בנו מסגרת להערכת ספקים אסטרטגיים, שמשלבת פריסת פונקציית האיכות (\en{QFD}) ו-\en{AHP}, ויישמו אותה בארגון ייצור בבריטניה \cite{dey2015}. לפיהם רוב המחקרים מעריכים ספקים לפי \textbf{מדדי תוצאה} (איכות, עמידה בלוח זמנים, עלות), ויש להוסיף \textbf{מדדים מקדימים}: ניהול הארגון, ניהול סיכונים ופרקטיקות סביבתיות וחברתיות. מדדי תוצאה דורשים היסטוריית עבודה, שאין בשלב טרום ההתקשרות, ולכן איכות מוערכת כאן לפי \textbf{אותות ציבוריים} שקרובים לרעיון המדדים המקדימים.

אות מרכזי הוא הסמכת מערכת ניהול האיכות, כמו \en{ISO 9001} או \en{AS9100} בענף התעופה. \en{Ayyadi} ו-\en{Oulhadj} סקרו 94 מחקרים מהשנים 2000 עד 2023 ומצאו שהשפעת הסמכת \en{ISO 9001} על ביצועי חברות ייצור \textbf{אינה אחידה}, ותלויה בסיבה לבחירה בהסמכה ובאופן השימוש במערכת \cite{ayyadi2024}.

\textbf{ביקורת.} הסמכה מראה שיש מערכת לניהול איכות, לא שהיא עובדת היטב. לכן ממד האיכות משלב כמה אותות: הסמכות בתוקף, לקוחות גדולים וחוזים ממשלתיים, פרסי ספקים, ואותות הפוכים כמו ריקולים והוראות תקינות של רשויות תעופה. התיקוף בודק אם ספקים מוכחים אכן מקבלים בממד זה ציון גבוה יותר.

% =====================================================================
\section{סיכון גיאופוליטי ורמת המדינה}\label{sec:lit-geo}

\en{Ho} ועמיתיו סקרו את ספרות ניהול סיכוני שרשרת האספקה בשנים 2003 עד 2013, ומהסקירה עולה שסיכון לספק נובע גם מגורמים חיצוניים, כמו המדינה שבה הוא פועל \cite{ho2015}. \en{Caldara} ו-\en{Iacoviello} בנו מדד סיכון גיאופוליטי (\en{GPR}) מספירת כתבות על מתחים, מלחמות וטרור \cite{caldara2022}, ובכך הראו שאפשר למדוד סיכון כזה כמותית מתוך טקסט חדשותי, אותו סוג מקור שהסוכנים יעבדו.

שני מחקרים שפיטים מראים שצריך לבדוק גם את \textbf{היחסים} בין מדינת הספק למדינת הלקוח. \en{Bai} ועמיתיו מצאו שלסיכון גיאופוליטי במדינת הספק יש השפעה כפולה על בסיס הספקים של חברות רב-לאומיות, והכיוון שלה תלוי בזיקה הפוליטית בין מדינת החברה למדינת הספק \cite{bai2025}. \en{Charpin} ו-\en{Cousineau} בדקו 2,858 חברות אמריקאיות שרוכשות מ-99 מדינות בשנים 2003 עד 2019, ומצאו שחברות מצמצמות ספקים במדינות שבהן המתיחות עולה \cite{charpin2025}. הממצאים תומכים בקטגוריה \textbf{יחסי המדינה עם ישראל}.

\textbf{ביקורת.} שני המחקרים בוחנים את בסיס הספקים של חברות גדולות לאורך שנים, ולא החלטה על ספק בודד. הם מצדיקים את הקטגוריה, אבל לא קובעים לה משקל; המשקל נקבע ב-\en{AHP}.

% =====================================================================
\section{חיזוי כשל פיננסי של ספקים}\label{sec:lit-fin}

עבודת היסוד היא מודל ה-\en{Z} של \en{Altman}, שמשקלל חמישה יחסים פיננסיים לציון אחד המבחין בין חברות שפשטו רגל לחברות יציבות \cite{altman1968}. אחד היחסים משתמש ב\textbf{שווי השוק} של ההון, ולכן אפשר לחשב אותו ישירות רק לחברה נסחרת. \en{Zhao}, \en{Ouenniche} ו-\en{De Smedt} סקרו את השיטות הסטטיסטיות ושיטות הבינה המלאכותית לחיזוי פשיטת רגל \cite{zhaoj2024}, והסקירה משמשת כאן כמפת השיטות האפשריות לממד הפיננסי.

\textbf{ביקורת: החברות הפרטיות.} כל השיטות האלה זקוקות לנתונים כספיים, ואילו לפי \en{RapidRatings} 22.1\% מהספקים הפרטיים בענף נמצאים בסיכון פיננסי גבוה, לעומת 15.9\% מהציבוריים \cite{rapidratings2026}. הסיכון גבוה דווקא במקום שבו הנתונים חסרים. לכן בחברה פרטית משמשים מדדים עקיפים, כמו אירועים שליליים, שינויי בעלות והסמכות שבוטלו, והדוח מציין כמה מהמידע נמצא.

% =====================================================================
\section{זיהוי אירועים שליליים מטקסט}\label{sec:lit-events}

סינון \textbf{חדשות שליליות} (\en{adverse media}) הוא רכיב קבוע בבדיקות "הכר את הלקוח" (\en{KYC}) ובמניעת הלבנת הון, ושם הוא נעשה בדרך כלל בחיפוש מילות מפתח שמייצר התרעות שווא רבות \cite{chernakov2025}. \en{Chernakov} ועמיתיו בנו סוכן מבוסס מודל שפה (\en{LLM}) עם אחזור מקורות (\en{RAG}), שמחפש ברשת ומחשב לכל נבדק מדד חדשות שליליות, ובדקו אותו על אנשי ציבור, על אנשים מרשימות פיקוח ומרשימות סנקציות ועל שמות "נקיים" \cite{chernakov2025}; זו טיוטה מוקדמת (\en{preprint}), והיא עוסקת באנשים ולא בספקים. \en{Zhang} ועמיתיו חילצו בעזרת מודל שפה 2,843 אירועי שיבוש מ-3,125 כתבות על שרשרת האספקה של רכבים חשמליים \cite{zhang2026}.

\textbf{ביקורת.} המחקרים עוסקים באנשים או בענף שלם, ואף אחד מהם אינו מתרגם אירועים להמלצה על \textbf{ספק בודד}. רובם משייכים כל אירוע לרשימה קבועה של קטגוריות (\textbf{טקסונומיה סגורה}), ולכן אירוע מסוג חדש נדחס לקטגוריה לא מתאימה. בכלי המוצע יש קטגוריה "אחר / חדש" שאדם מאשר, והחומרה נקבעת לפי הפגיעה ביכולת לספק ולא לפי סוג האירוע: מתקפת סייבר שעצרה ייצור היא כשל אספקה, ומתקפה שלא עצרה אותו היא אות סיכון.

% =====================================================================
\section{אבטחת סייבר של ספקים}\label{sec:lit-cyber}

\en{ENISA} מצאה שב-66\% מ-24 מתקפות על שרשראות אספקה הספקים לא ידעו או לא דיווחו איך נפרצו \cite{enisa2021}, ובענף התעופה מתקפת כופרה (\en{ransomware}) על \en{ASCO} עצרה ייצור בכמה מדינות \cite{certeu2019}. המסגרת הרשמית בארה"ב היא \en{NIST SP 800-161 Rev.~1} \cite{nist2022}, ומשרד ההגנה האמריקאי מחייב את ספקיו בהסמכת \en{CMMC} \cite{cmmc2024}, שהיא גם \textbf{אות ציבורי} שאפשר לבדוק בלי לפנות לספק. \en{Gani} ועמיתיו מצאו, בסקר של 130 יצרנים במלזיה, ש\textbf{נראוּת} של שרשרת האספקה קשורה לביצועי אבטחת הסייבר שלה \cite{gani2023}. \en{Slapni{\latinfallback č}ar} ועמיתיו ראיינו 33 מומחים וסקרו 53 נוספים, ובנו תיאוריה של תהליך הערכת סיכוני הסייבר של ספקים \cite{slapnicar2025}.

\textbf{ביקורת.} רוב התהליכים נשענים על \textbf{שאלונים שהספק ממלא בעצמו} ועל גישה למערכות שלו, ואין גישה כזו לפני ההתקשרות; ממצא \en{ENISA} מראה שגם ספק שמשיב לא תמיד יודע מה קרה אצלו. לכן הכלי בודק סייבר \textbf{מבחוץ ובאופן פסיבי בלבד}: הסמכות ודיווחים פומביים על פריצות.

% =====================================================================
\section{מודלי שפה ומערכות מרובות סוכנים בשרשרת אספקה}\label{sec:lit-agents}

\en{Jackson} ועמיתיו הציעו מסגרת רעיונית, שלא נבדקה על נתונים, לשימוש בבינה מלאכותית יוצרת ב-13 תחומי החלטה בשרשרת אספקה \cite{jackson2024}. \en{AlMahri}, \en{Xu} ו-\en{Brintrup} בנו מערכת של \textbf{שבעה סוכנים} שמזהים אותות שיבוש בחדשות, ממפים אותם לרשת ספקים וממליצים על צעדים \cite{almahri2026}. לפיהם המערכת הגיעה ל-\en{F1} של 0.962 עד 0.991, וניתוח ארך 3.83 דקות ועלה כ-0.08 דולר. זו טיוטה שלא עברה שיפוט, שנבדקה על 30 תרחישים \textbf{מסונתזים} בתעשיית הרכב, והיא \textbf{מנטרת ספקים קיימים} ולא בודקת ספק חדש.

\textbf{הסיכון של מודלי שפה.} \en{Huang} ועמיתיו סקרו את תופעת ה\textbf{הזיה} (\en{hallucination}): טקסט שנשמע אמין אבל אינו נכון, שגם אחזור מקורות אינו מבטל לגמרי \cite{huang2025}. בהמלצה על ספק ביטחוני טענה מומצאת היא כשל חמור, ולכן כל טענה מקושרת למקור (\en{REQ-NFR-01}), וטענה בלי מקור אינה נכנסת לציון.

% =====================================================================
\section{מטריצת הספרות}\label{sec:lit-matrix}

טבלה~\ref{tab:litmatrix} מרכזת את המחקרים העיקריים מול הנושאים שהבעיה דורשת.

{\small
\begin{longtable}{|>{\raggedright\arraybackslash}p{2.4cm}|>{\centering\arraybackslash}p{1.12cm}|>{\centering\arraybackslash}p{1.12cm}|>{\centering\arraybackslash}p{1.12cm}|>{\centering\arraybackslash}p{1.12cm}|>{\centering\arraybackslash}p{1.12cm}|>{\centering\arraybackslash}p{1.12cm}|>{\centering\arraybackslash}p{1.12cm}|>{\centering\arraybackslash}p{1.12cm}|}
\caption{מטריצת הספרות: השורות הן מחקרים, והעמודות הן נושאים שהבעיה דורשת}\label{tab:litmatrix}\\
\hline
\textbf{מחקר} & \textbf{טקסט חופשי ומודל שפה} & \textbf{המלצה על ספק בודד} & \textbf{שלב טרום התקש\-רות} & \textbf{איכות ואמינות הספק} & \textbf{חברות פרטיות וכיסוי מידע} & \textbf{סייבר} & \textbf{גיאו\-פוליטי} & \textbf{תיקוף על כשלים אמיתיים} \\ \hline
\endfirsthead
\hline
\textbf{מחקר} & \textbf{טקסט חופשי ומודל שפה} & \textbf{המלצה על ספק בודד} & \textbf{שלב טרום התקש\-רות} & \textbf{איכות ואמינות הספק} & \textbf{חברות פרטיות וכיסוי מידע} & \textbf{סייבר} & \textbf{גיאו\-פוליטי} & \textbf{תיקוף על כשלים אמיתיים} \\ \hline
\endhead
\hline
\endlastfoot
\en{Ho} ועמיתיו \cite{ho2010} & --- & נבדק & נבדק & חלקי & --- & --- & --- & --- \\ \hline
\en{Dey} ועמיתיו \cite{dey2015} & --- & נבדק & --- & נבדק & --- & --- & --- & --- \\ \hline
\en{Altman} \cite{altman1968} & --- & חלקי & --- & --- & --- & --- & --- & נבדק \\ \hline
\en{Bai} ועמיתיו \cite{bai2025} & --- & --- & --- & --- & --- & --- & נבדק & --- \\ \hline
\en{Slapni{\latinfallback č}ar} ועמיתיו \cite{slapnicar2025} & --- & חלקי & --- & --- & --- & נבדק & --- & --- \\ \hline
\en{Chernakov} ועמיתיו \cite{chernakov2025} & נבדק & חלקי & חלקי & --- & --- & --- & --- & חלקי \\ \hline
\en{AlMahri} ועמיתיו \cite{almahri2026} & נבדק & --- & --- & --- & --- & --- & חלקי & --- \\ \hline
\textbf{מאמר נלווה:} \en{Zhao} ועמיתיו \cite{zhao2025} & נבדק & --- & --- & --- & --- & --- & חלקי & --- \\ \hline
\textbf{מאמר הדגל:} \en{Sadeghi} ועמיתיו \cite{sadeghi2026} & חלקי & --- & נבדק & נבדק & --- & --- & חלקי & --- \\ \hline
\textbf{הפרויקט המוצע} & ייבדק & ייבדק & ייבדק & ייבדק & \textbf{ייבדק: זה הפער} & ייבדק & ייבדק & ייבדק \\
\end{longtable}
}

מחקרי \en{MCDM} ואיכות ספקים מגיעים להמלצה ומעריכים חוזקות, אבל מניחים מידע מלא והיסטוריית ביצוע. מחקרי מודלי השפה מחפשים רק סיכונים. מאמר הדגל מחבר בין העולמות, אבל אינו אוסף מידע ציבורי ואינו מטפל במידע חסר. אף מחקר אינו מדווח כמה מהמידע על הספק נמצא.

% =====================================================================
\section{מאמר הדגל}\label{sec:flagship}

\textbf{הבחירה.} מאמר הדגל הוא המאמר של \en{Sadeghi}, \en{Dehghani} ו-\en{Kabadayi} (2026) בכתב העת \en{Annals of Operations Research} \cite{sadeghi2026}. הוא \textbf{עוסק באותה בעיה}: בחירת ספקים לפי כמה קריטריונים בעזרת מודל שפה ושקלול ב-\en{AHP}, כשאיכות ואמינות הם אחד הממדים; הוא \textbf{עדכני ופורסם בכתב עת שפיט מוביל}; ו\textbf{אפשר לשחזר אותו}, כי השלבים, משוואות ה-\en{AHP} וסף העקביות מתוארים במלואם, והתוצאות הושוו לשיפוט מומחים במבחן סטטיסטי.

\textbf{המודל.} המסגרת \en{AIDM} (\en{AI-driven Decision-Making}) מחליפה במודל שפה את פאנל המומחים שממלא את ה-\en{AHP}, בארבעה שלבים: (1) המודל יוצר חמישה מומחים וירטואליים עם רקע שונה; (2) כל מומחה מציע קריטריונים, והחשובים ביותר מקובצים להיררכיה; (3) המומחים משווים בזוגות, והשוואות לא עקביות חוזרות עד חמש פעמים; (4) המשקלות מושווים למשקלות שקבעו מומחי רכש אמיתיים. ההיררכיה שהתקבלה כוללת שלושה ממדים: \textbf{איכות} (איכות מוצר, אמינות, תקשורת, גמישות ומשוב לקוחות), \textbf{עלות} (מחיר, זמן אספקה, עדכוני מחיר ותנאי חוזה) ו\textbf{מאפייני הספק} (השפעה סביבתית, נגישות, \textbf{מדינת המקור}, ממליצים, מערכות ניהול וניסיון גלובלי).

\textbf{המשוואות השולטות.} כל מעריך ממלא מטריצת השוואות $A$ בגודל $n\times n$, שבה $a_{ij}$ אומר כמה קריטריון $i$ חשוב מקריטריון $j$. המשקלות הם הווקטור העצמי של הערך העצמי הגדול ביותר $\lambda_{\max}$, והעקביות נמדדת ביחס $\mathrm{CR}$, כש-$\mathrm{RI}$ הוא ערך ייחוס שתלוי בגודל המטריצה \cite{saaty1990}:
\begin{equation}\label{eq:ahp}
A\,\mathbf{w}=\lambda_{\max}\,\mathbf{w},\qquad \sum_{i=1}^{n} w_i = 1
\end{equation}
\begin{equation}\label{eq:cr}
\mathrm{CI}=\frac{\lambda_{\max}-n}{n-1},\qquad
\mathrm{CR}=\frac{\mathrm{CI}}{\mathrm{RI}}\le 0.10
\end{equation}
סף העקביות במאמר הוא 0.10, והוא מאומץ גם כאן. כדי להשוות בין משקלות המומחים האנושיים $w^{H}$ לווירטואליים $w^{AI}$, חושב לכל קריטריון $|w^{H}_i-w^{AI}_i|$, וההבדל נבדק במבחן \en{Wilcoxon} לזוגות \cite{sadeghi2026}.

\textbf{הניסוי והתוצאות.} המקרה הוא בחירת ספקים בחברה לייצור מכונות כבדות; המומחים הווירטואליים נבנו עם \en{\mbox{ChatGPT-4}}, וקבוצת ההשוואה כללה 14 אנשי רכש (8 ברכש המקומי ו-6 בבינלאומי) \cite{sadeghi2026}. שתי הקבוצות עמדו בסף העקביות, והאנושיים היו עקביים יותר ($\mathrm{CR}=0.030$ ו-$0.060$); 85\% מהמומחים הווירטואליים הגיעו ליחס תקין. מבחן \en{Wilcoxon} לא מצא הבדל מובהק ($\alpha=0.05$) בין קבוצות המשקלות, אבל ברמת הקריטריון היו הבדלים עקביים (טבלה~\ref{tab:sadeghi-results}).

{\small
\begin{longtable}{|>{\raggedright\arraybackslash}p{3.6cm}|c|c|c|c|}
\caption{דוגמאות למשקלות במאמר הדגל: מומחים אנושיים מול וירטואליים \cite{sadeghi2026}}\label{tab:sadeghi-results}\\
\hline
\textbf{קריטריון} & \textbf{מקומי: אנושי} & \textbf{מקומי: וירטואלי} & \textbf{בינלאומי: אנושי} & \textbf{בינלאומי: וירטואלי} \\ \hline
\endfirsthead
\hline
\textbf{קריטריון} & \textbf{מקומי: אנושי} & \textbf{מקומי: וירטואלי} & \textbf{בינלאומי: אנושי} & \textbf{בינלאומי: וירטואלי} \\ \hline
\endhead
אמינות & 0.114 & 0.064 & 0.139 & 0.075 \\ \hline
עדכוני מחיר & 0.099 & 0.031 & 0.107 & 0.033 \\ \hline
מדינת המקור & 0.029 & 0.065 & 0.024 & 0.066 \\ \hline
מודעות סביבתית & 0.024 & 0.050 & 0.016 & 0.051 \\* \hline
\end{longtable}
}

אנשי הרכש נתנו משקל גבוה פי שניים בערך ל\textbf{אמינות}, ומודל השפה משקל גבוה פי שניים בערך ל\textbf{מדינת המקור}. אף שהמבחן לא מצא הבדל ברמת כלל המשקלות, ההבדלים חשובים דווקא בגלל הקריטריונים שבהם הופיעו: אלה שני הנושאים הרגישים ביותר לתעשייה ביטחונית. לכן במתודולוגיה המוצעת המשקלות נקבעים בשיפוט אנושי המנומק בספרות, ולא על ידי הסוכנים. מהמאמר נלקחים המבנה ההיררכי, משוואות ה-\en{AHP}, סף העקביות וממדי האיכות והאמינות; הרעיון של מודל שפה כמעריך משקלות נדחה.

\textbf{המגבלות.} המחברים מציינים שמודל השפה עלול להמציא קריטריונים, ושאין ערובה שהרשימה מלאה \cite{sadeghi2026}. בנוסף זוהו שש מגבלות: (1) המאמר \textbf{עוצר במשקלות} ואינו מעריך אף ספק בפועל; (2) המסגרת \textbf{אינה אוספת מידע} על הספק; (3) \textbf{אין טיפול במידע חסר}; (4) \textbf{סיכון הוא משקל בלבד}, כך שספק ממדינה מסוכנת יכול "לפצות" במחיר נמוך, ואין שער חסימה, אירועים חריגים או סייבר; (5) התיקוף הוא \textbf{מול מומחים ולא מול המציאות}; (6) \textbf{חברה אחת וענף אחד}.

% =====================================================================
\section{המאמר הנלווה: זיהוי סיכונים מחדשות}\label{sec:companion}

מאמר הדגל עוסק בעיקר בפרופיל החוזקות ובמשקלות. לחלק דגלי הסיכון משמש מאמר נלווה: \en{Zhao} ועמיתיו (2025), \en{Service Oriented Computing and Applications} \cite{zhao2025}, שבודק אם מודלי שפה מזהים נכון סיכוני ספקים מתוך כתבות אמיתיות, כלומר את משימת סוכני הסיכון.

המסגרת \en{LARD-SC} אוספת כתבות על כל ספק מ-\en{Google News}; מודל שפה מסכם כל כתבה, מזהה סיכון ומעריך את סבירותו והשפעתו; וכל סיכון ממופה לטקסונומיה של \en{Cambridge} לסיכונים עסקיים. מתוך 676 כתבות על 188 ספקים בשכבה הראשונה של \en{Apple} (2020--2024) נדגמו 120, חמישה מודלים קיבלו אותן עם אותו פרומפט, ושני זוגות מומחים בדקו כל סיכון. חושבו שיעור הסיכונים שאושרו (\en{RVR}), שאושרו חלקית (\en{PRR}), הזיהויים השגויים (\en{FIR}) ומדד ביצוע יחסי (\en{RPI}) (טבלה~\ref{tab:zhao-results}).

{\small
\begin{longtable}{|l|c|c|c|c|}
\caption{תוצאות המאמר הנלווה: ממוצע שני זוגות המומחים \cite{zhao2025}}\label{tab:zhao-results}\\
\hline
\textbf{מודל} & \en{RVR} & \en{PRR} & \en{FIR} & \en{RPI} \\ \hline
\endfirsthead
\hline
\textbf{מודל} & \en{RVR} & \en{PRR} & \en{FIR} & \en{RPI} \\ \hline
\endhead
\hline
\endlastfoot
\en{GPT-4o} & 0.8875 & 0.0917 & 0.0208 & 0.8484 \\ \hline
\en{GPT-4o mini} & 0.8750 & 0.1042 & 0.0208 & 0.8334 \\ \hline
\en{Claude 3.5 Sonnet} & 0.8708 & 0.1042 & 0.0250 & 0.8250 \\ \hline
\en{Claude 3.5 Haiku} & 0.8250 & 0.1542 & 0.0208 & 0.7734 \\ \hline
\en{GPT-3.5 turbo} & 0.4625 & 0.2417 & 0.2958 & 0.1184 \\
\end{longtable}
}

במאמר מוגדר \en{$\mathrm{RPI}=w_1\mathrm{RVR}-w_2\mathrm{PRR}-w_3\mathrm{FIR}$} עם המשקלות \en{$w_1=1,\ w_2=0.5,\ w_3=1$}, אבל ערכי ה-\en{RPI} בטבלה 3 של המאמר, שמוצגים כאן, מתקבלים רק עם \en{$w_2=0.2$}; למשל ל-\en{GPT-4o}: \en{$0.8875-0.2\cdot0.0917-0.0208=0.8484$}. דירוג המודלים זהה בשתי הדרכים. ארבעת המודלים החדשים הגיעו לשיעור אימות של 0.82 עד 0.89 עם פחות מ-3\% זיהויים שגויים, והישן נכשל בכמעט 30\% מהמקרים: \textbf{בחירת המודל משנה את התוצאה מהותית}, ולכן התיקוף ייעשה בשני מודלים לפחות (\en{T4.1}, פרק~\ref{ch:workplan}).

\textbf{המגבלות.} מעבר להשפעת ניסוח הפרומפט ולתלות בזמינות המידע, שהמחברים מציינים \cite{zhao2025}: (1) \textbf{רק דיוק, בלי כיסוי}: אין מדד לסיכונים שהוחמצו, ובספק ביטחוני החמצה חמורה מהתרעת שווא; (2) המומחים בדקו אם הסיכון \textbf{סביר}, ולא אם התממש; (3) הספקים \textbf{גלויים}, וכשאין כתבות, היעדר חדשות רעות עלול להתפרש כהיעדר סיכון; (4) הניתוח ברמת \textbf{הכתבה}, בלי ציון ספק ובלי שערי חסימה; (5) \textbf{רק סיכונים}, בלי חוזקות.

% =====================================================================
\section{פער הידע}\label{sec:gap}

\begin{quote}
\textbf{עד כה לא נבדק} אם מערכת מבוססת מודלי שפה, שנשענת על מידע ציבורי בלבד, יכולה להפיק \textbf{המלצת התקשרות על ספק בודד} בשלב טרום ההתקשרות, להעריך בה גם את \textbf{איכות הספק ואמינותו} ולא רק סיכונים, ולדווח במפורש \textbf{כמה מהמידע הנדרש נמצא}, כשהספק הוא חברה פרטית שעליה מעט מידע, ולתקף את ההמלצה מול \textbf{מקרי כשל אמיתיים}.
\end{quote}

הפרויקט עונה על המגבלות שתוארו: הוא ממשיך ממשקלות להערכה מלאה של ספק מסוים; הסוכנים אוספים מידע ציבורי וכל טענה מקושרת למקור; סנקציות, בקרת יצוא ומדינה עוינת הם שער פסילה ולא קריטריון בשקלול; ציון הכיסוי מונע מהיעדר כתבות להתפרש כהיעדר סיכון; בחברה פרטית משמשים מדדים עקיפים; והתיקוף נעשה מול מקרי כשל ציבוריים ומול ספקים מוכחים, כך שנמדדות גם החמצות.
